\section{Post-Pause Experiment Blueprint}

This blueprint specifies the concrete execution order for the first post-pause campaign. It turns the evaluation plan into a staged workflow with stop/go criteria, required outputs, and artifact expectations.

\subsection{Stage 0: Environment Freeze and Sanity Checks}

Before any measurement run, freeze implementation commit, paper commit, model artifact identity, backend build metadata, and runtime configuration templates. Run a small deterministic sanity suite to confirm that encode/decode behavior matches expected fail-closed semantics under known-good and known-bad inputs.

Stage 0 exits only when sanity checks pass and metadata capture is validated.

\subsection{Stage 1: Recoverability Baseline}

Run the core condition grid with default conservative policy. Primary outputs are decode success rate, failure-class histogram, and retry distribution by condition.

A practical stop criterion for this stage is that every condition has enough samples to estimate success rate with stable confidence intervals and no missing failure taxonomy fields.

\subsection{Stage 2: Runtime and Efficiency Profiling}

Using the same sample grid, measure encode/decode latency distributions, throughput, packet-token efficiency, and total-token efficiency. Report median, p90, and p99 summaries with raw distributions in released artifacts.

This stage should also break down overhead by segment type (bootstrap vs body) to identify where deterministic safeguards have the largest runtime cost.

\subsection{Stage 3: Approximation-Term Instrumentation}

Enable step-level logging for $\alpha_t$, $N_t$, $F_{\mathrm{tot}}$, and $\beta_t$.
Compute per-message summaries and empirical diagnostics based on
\begin{equation}
\widehat{B}_{\mathrm{TV}}=\sum_t\left(\alpha_t+\frac{N_t}{2F_{\mathrm{tot}}}+\beta_t\right).
\end{equation}

The main output is a decomposition of which term dominates in which operating regions, not a single aggregate score.

\subsection{Stage 4: Ablation Campaign}

Execute the ablation matrix one component at a time:
\begin{itemize}
\item stability filter,
\item fingerprint check,
\item retry policy,
\item quantization budget,
\item truncation policy,
\item tail policy.
\end{itemize}

For each ablation, report deltas in recoverability, efficiency, failure composition, and approximation-term summaries. This stage is the main evidence path for novelty-through-outcome arguments.

\subsection{Stage 5: Detector-Facing Evaluation}

Evaluate a detector suite under single-message and multi-message accumulation settings. Detector outputs should be paired with approximation and traffic metrics from previous stages so shifts can be interpreted rather than only reported.

Multi-message results must be shown as curves over aggregation size rather than a single pooled point.

\subsection{Stage 6: Traffic-Side Analysis}

Using retry and timing logs, compare stego and non-stego baselines on output-length and latency distributions. Report unconditional and retry-conditioned distributions to separate intrinsic pipeline effects from retry-mix effects.

If this stage reveals strong side-channel separation, mitigation tuning should be prioritized before any detectability claims are strengthened.

\subsection{Stage 7: Artifact Packaging and Internal Reproduction}

For each figure/table, package raw data, aggregation scripts, plotting scripts, and a README with exact commands. A second local reproduction pass should rebuild all tables and figures from raw logs to validate artifact completeness.

Only after this stage should claim language in the main text be upgraded.

\subsection{Result Templates for Reporting}

To keep reporting consistent, each result block should include:
\begin{itemize}
\item condition definition,
\item sample size,
\item primary metrics with uncertainty,
\item failure-class decomposition,
\item relevant approximation/traffic diagnostics,
\item artifact path and commit hashes.
\end{itemize}

\subsection{Escalation Rules for Weak Outcomes}

The blueprint should treat weak outcomes as first-class results.
If recoverability is unstable, prioritize policy hardening before detector experiments.
If efficiency is low but recoverability is strong, prioritize quantization and truncation tuning with failure safeguards unchanged.
If detector or traffic risks rise quickly with payload pressure, prioritize conservative operating envelopes and explicit caveats.

\subsection{Expected Deliverable of the First Post-Pause Cycle}

The first cycle should produce a fully auditable evidence package: measured recoverability, runtime and efficiency profiles, approximation-term decomposition, ablation deltas, detector-facing results, traffic-side analysis, and reproducible artifacts. This package is the minimum basis for replacing design-time statements with evidence-backed claims.
